\subsection{Manejadores \texttt{async void}}

\subsubsection{Regla: un manejador \texttt{async void} no deja escapar excepciones}

\textbf{Regla:} \texttt{async void} se usa únicamente en manejadores de eventos. Todo su cuerpo queda dentro de un \texttt{try} que captura cada excepción que la operación puede lanzar y la maneja según la Sección 11.1.

\textbf{Descripción:} Un método \texttt{async void} no devuelve una \texttt{Task} donde guardar la excepción: la lanza directamente sobre el contexto de sincronización activo, ningún \texttt{catch} del llamador la recibe y el proceso termina (Cleary, 2013).

\textbf{Código correcto:}
\begin{lstlisting}[style=csharpstyle]
private async void HandleRollClick(object? sender, EventArgs e)
{
    try
    {
        await _dice.RollAsync();
    }
    catch (TimeoutException ex)
    {
        _logger.Warn("The dice roll did not finish in time.", ex);
        ShowRetryNotice();
    }
}
\end{lstlisting}

\textbf{Código incorrecto:}
\begin{lstlisting}[style=csharpstyle]
private async void HandleRollClick(object? sender, EventArgs e)
{
    await _dice.RollAsync();
}
\end{lstlisting}
